\section{Introduction}
\label{sec:intro}

%%% Paragraph 1: Impact statement + Territory %%%
VLA models are being deployed in physical environments where they
manipulate objects alongside
humans~\citep{rt2,openvla,pi0}, introducing \emph{physical
irreversibility}: a wrong robotic action produces consequences that
no software rollback can undo~\citep{amodei2016concrete}.
Meanwhile, VLA weights are increasingly updated through cloud-edge
continuous learning pipelines, where models fine-tuned on
edge-collected demonstrations are periodically pushed to deployed
systems---a practice that makes post-deployment backdoor injection
both feasible and difficult to audit.
Visual backdoor attacks~\citep{badvla,goba,tabvla,attackvla}
exploit this exposure: an adversary poisons training demonstrations to
bind a visual trigger to a dangerous action sequence, reporting
attack success rates above 97\%~\citep{goba,badvla} while leaving
clean-input task performance---the signal on which standard
evaluations rely---essentially unchanged.
Pre-deployment defenses presuppose centralized training access
that continuous cloud-edge updates no longer provide, and existing
inference-time monitors either target \emph{accidental} failures
or intervene in the visual input space before encoding, leaving
post-encoding hidden-state and trajectory-level signals largely
untapped for backdoor defense~(\S\ref{sec:related}).
We propose \textbf{Security Cuff}, an inference-time defense for
deployed VLA systems that detects and intercepts backdoor-triggered
actions before they cause irreversible physical harm.

%%% Paragraph 2: Method overview + key insight %%%
Security Cuff is grounded in a \emph{physical consequence
invariant}: for the class of VLA backdoor attacks we target---
those whose payload is a physically dangerous action---a
successful trigger must ultimately manifest as a harmful
trajectory, even if the trigger itself is hidden.
We exploit this through a dual-process
architecture~\citep{kahneman2011thinking}: a fast statistical
screen (System~1) paired with a slow consequence-aware verifier
(System~2).
The \emph{fast layer}, with a target latency of ${\leq}80\,$ms per
step, runs an Augmented SAFE probe~\citep{safe} that detects
hidden-state distributional shifts from clean execution---a signal
distinct from the output-level uncertainty used by OOD monitors
~\citep{mahal,ensemble_ood}---alongside a Temporal Trajectory
Divergence~(TTD) monitor over short-to-medium-horizon
\emph{observed} trajectory statistics to flag anomalous execution
drift~\citep{attackvla}.
The \emph{slow layer}, with a target latency of ${\leq}300\,$ms,
rolls out a world model~\citep{ha2018world} and scores the
\emph{predicted} long-horizon trajectory against physical safety
criteria, directly operationalizing the physical consequence
invariant as a runtime check.

%%% Paragraph 3: Evidence preview %%%
% TODO: fill in after experiments
We evaluate Security Cuff on \texttt{[benchmark]} against
\verb|[N]| state-of-the-art VLA backdoor attacks spanning
physical-object triggers~\citep{goba}, visual
patches~\citep{tabvla}, pixel
perturbations~\citep{badvla}, and stealth trajectory
deviations~\citep{silentdrift}.
We additionally report auxiliary out-of-distribution failure
experiments to probe whether the detector transfers beyond the
adversarial setting.
Security~Cuff reduces the catastrophic failure rate to
\verb|[X%]| while maintaining task success within
\verb|[Y%]| of the clean baseline at \verb|[Z] ms|
overhead per step.
Ablations confirm both layers are necessary: the fast layer alone
misses stealth attacks; the slow layer alone exceeds real-time
latency budgets.

%%% Paragraph 4: Contributions %%%
\paragraph{Contributions.}
\begin{enumerate}
\item \textbf{A trajectory-level formalization of physically
dangerous actions under VLA backdoor attacks.}
We formalize \emph{physically dangerous actions} induced by
backdoor triggers in deployed VLA systems and articulate the
\emph{physical consequence invariant}: any backdoor whose payload
is a dangerous action must ultimately surface as a harmful
trajectory, regardless of trigger form. This reframes
trigger-agnostic backdoor interception as a trajectory-safety
problem and gives the paper its primary deployment-time threat
model; auxiliary out-of-distribution failures are used only for
secondary robustness experiments.

\item \textbf{A consequence-aware verifier based on world-model
rollout.}
We instantiate a \emph{consequence-aware verifier} by rolling out
an auxiliary dynamics model from candidate prefixes and scoring the
predicted trajectories against physical safety criteria. Unlike
trigger-reconstruction or uncertainty-only monitors, this verifier
reasons directly about downstream harm and requires neither trigger
priors nor access to training data or weights.

\item \textbf{A fast--slow runtime system for multi-timescale
backdoor coverage.}
We design Security~Cuff as a plug-and-play fast--slow runtime
system wrapping an unmodified VLA policy. A low-latency hidden-state
screen runs on every step, and only suspicious prefixes are
escalated to the consequence-aware verifier. The two layers cover
both abruptly triggered and stealthy long-horizon backdoors within
real-time latency budgets and without retraining or weight access.
\end{enumerate}
